\section{Вычислительный юмор как исследовательская проблема}

Юмор является одной из наиболее трудных задач для обработки естественного языка, поскольку успешная шутка определяется не только грамматической правильностью или тематической релевантностью. Она обычно требует нарушения ожиданий, использования общего знания, выбора подходящей формы, учета аудитории и точного соотношения между подготовкой и неожиданным завершением. Поэтому генерация юмора занимает промежуточное положение между задачами языкового моделирования, творческого письма и прагматического анализа коммуникации.

Ранние исследования вычислительного юмора уже подчёркивали, что юмор плохо сводится к одному формальному правилу. \citet{ritchie2001current} описывает вычислительный юмор как область, в которой частные механизмы могут быть формализованы, но общая способность понимать и порождать юмор требует знания о языке, мире и намерениях участников коммуникации. Обзор \citet{amin2020survey} показывает, что классические системы генерации юмора часто строились вокруг узких механизмов: шаблонов, лексических ресурсов, фонетических замен, каламбуров или заранее заданных форм шутки. Такие системы были интерпретируемыми и позволяли контролировать юмористический механизм, но плохо переносились на открытые темы и разнообразные запросы.

Современные обзоры приходят к сходному выводу уже в контексте больших языковых моделей. \citet{Loakman2025} отмечают, что генерация и объяснение юмора остаются недостаточно изученными по сравнению с более стандартными задачами обработки естественного языка, особенно если выходить за пределы каламбуров. \citet{lemmens2026computational} также подчёркивают, что несмотря на рост возможностей нейросетевых моделей, юмор остаётся областью с неоднозначными результатами: модели способны воспроизводить внешние признаки юмористического текста, но это не означает устойчивого понимания юмористических механизмов.

Каламбуры являются удобным примером этой проблемы. С одной стороны, они дают относительно чёткий объект анализа: в них можно изучать неоднозначность, фонетическое сходство и переключение смысла. Так, \citet{kao2016computational} формализуют юмор в каламбурах через сочетание неоднозначности и отличительности, показывая, что эти факторы связаны с человеческими оценками смешного. С другой стороны, успех на каламбурах не переносится автоматически на более широкие формы юмора: короткие нарративные шутки, тематические шутки, ситуационный юмор или социально обусловленные ответы. Современные исследования больших языковых моделей подтверждают это ограничение. \citet{xu2024good} показывают, что модели часто демонстрируют «ленивую» генерацию каламбуров: вместо компактного совмещения двух значений они явно называют оба связанных понятия или воспроизводят поверхностный шаблон игры слов. \citet{cocchieri2025you} также показывают, что даже сильные языковые модели испытывают трудности с обобщением в задачах понимания и генерации юмора, когда требуется не просто распознать знакомую форму, а применить юмористический механизм в новой ситуации.

Для данной работы из этого следует два вывода. Во-первых, генерацию юмора нельзя рассматривать как обычную задачу продолжения текста: беглость и тематическая релевантность являются необходимыми, но недостаточными условиями. Во-вторых, юмор является задачей с множеством допустимых ответов. Для одного и того же запроса может существовать большое число хороших шуток, основанных на разных ассоциациях и разных механизмах неожиданности. Это делает постановку с единственным эталонным ответом методологически слабой для обучения и оценки юмористической генерации.

\section{Генерация юмора большими языковыми моделями}

Появление больших языковых моделей изменило характер исследований генерации юмора. В отличие от шаблонных систем, такие модели способны отвечать на широкий спектр запросов и порождать тексты, которые по форме похожи на человеческие шутки. Однако литература показывает, что это расширение охвата не устраняет центральные трудности области.

Крупные корпуса юмора, такие как rJokes, важны как источники разнообразных реальных шуток \citep{weller2020rjokes}. Они позволяют изучать тематическое разнообразие, повторяющиеся формы и статистические свойства юмористических текстов. Однако такие корпуса обычно не строятся как наборы данных для многореференсной генерации. В них нет явного соответствия между одним запросом и множеством допустимых ответов. Поэтому прямое дообучение модели на корпусе шуток может научить её стилю и распространённым формам, но не обязательно научит выбирать неожиданный и уместный юмористический ход для нового запроса.

Ограничения современных моделей хорошо видны в задачах, где существует сильная человеческая оценка. \citet{Zhang2024HumorAI} представляют крупный ресурс для генерации подписей к карикатурам, содержащий более 250 миллионов человеческих оценок и более 2.2 миллиона подписей. Эта работа особенно важна, поскольку рассматривает юмор как творческую задачу предпочтений, а не как бинарную классификацию. Авторы показывают, что даже сильные модели и методы дообучения по предпочтениям остаются ограниченными на творческой задаче: лучшие системы всё ещё уступают сильным человеческим участникам.

Похожие выводы появляются и в более специализированных исследованиях. \citet{loakman2025comparing} анализируют понимание юмора от традиционных каламбуров до тематических шуток и показывают, что современные модели не одинаково хорошо справляются с разными типами юмора. \citet{Sakabe2025} рассматривают японскую импровизационную форму Oogiri и оценивают ответы по нескольким измерениям: новизне, ясности, релевантности, интеллектуальности, эмпатии и общей смешности. Их результаты показывают, что языковые модели могут достигать уровня между низкими и средними человеческими результатами, но заметно отличаются от людей по тому, какие качества ответа они считают важными.

Эти работы показывают, что прогресс больших языковых моделей не делает задачу генерации юмора решённой. Скорее, он меняет исследовательский вопрос. Если раньше основная трудность состояла в том, чтобы вообще получить грамматически и структурно похожий на шутку текст, то теперь важно отличать поверхностную имитацию от устойчивого творческого обобщения. Для данной работы это особенно существенно: оптимизируемая модель не должна просто запоминать распространённые шутки или воспроизводить корпусные шаблоны, а должна использовать эталонные шутки как слабый сигнал для обучения распределения творческих траекторий.

\section{Проблема оценки юмора}

Оценка юмора является отдельной методологической трудностью. В отличие от задач с объективно проверяемым ответом, смешность зависит от аудитории, контекста, культурных ожиданий, формы подачи и конкретного измерения качества. Поэтому простая формулировка «ответ смешной или нет» часто оказывается слишком грубой.

Корпус rJokes уже демонстрирует эту проблему: реальные шутки разнообразны по теме, форме и качеству, а их восприятие зависит от сообщества, в котором они были опубликованы \citep{weller2020rjokes}. Более поздние работы делают этот вывод сильнее. \citet{Sakabe2025} показывают, что оценка юмора должна учитывать несколько измерений, а не только общую смешность. Их анализ также показывает расхождение между человеческими и модельными критериями оценки: люди и языковые модели могут придавать разный вес новизне, эмпатии и релевантности.

\citet{winters2025evaluating} предлагают другой важный взгляд, сравнивая генерацию юмора в условиях, близких к импровизационной комедии. Такая постановка подчёркивает, что юмор не является только свойством изолированного текста. На оценку влияет ситуация, скорость ответа и ожидания аудитории. Следовательно, результаты статической текстовой оценки не следует напрямую отождествлять с качеством юмора в живом взаимодействии.

В то же время полностью отказаться от автоматической оценки в вычислительных экспериментах трудно. Человеческая оценка дорога, шумна и плохо масштабируется. Поэтому современные работы используют языковые модели как вспомогательных оценщиков или строят более структурированные рубрики. \citet{Hashemi2024} предлагают многомерный и калиброванный подход к автоматической оценке текстов, показывая, что оценивание с помощью языковых моделей требует явной структуры и контроля. \citet{Gunjal2025} рассматривают рубрики как источник награды для обучения за пределами строго проверяемых доменов. Эти подходы важны для данной работы, поскольку юмор также относится к областям, где внешняя награда является неполной и шумной.

Тем не менее модель-оценщик не является прозрачной заменой человеческого восприятия. В данной работе автоматическая попарная оценка используется как воспроизводимый инструмент промежуточного сравнения моделей, а не как окончательное измерение человеческой смешности. Это ограничение принципиально: даже если модель выигрывает в попарной оценке, такой результат следует интерпретировать как улучшение относительно выбранного протокола оценки, а не как доказательство универсального юмористического качества.

\section{Рассуждение и творческий поиск в генерации юмора}

Один из важных сдвигов последних лет состоит в переходе от мгновенной генерации ответа к явному моделированию промежуточного рассуждения. В широком контексте обработки естественного языка \citet{Wei2022} показывают, что подсказки с цепочкой рассуждения могут существенно менять поведение больших языковых моделей на задачах, требующих многошагового вывода. Для юмора этот результат нельзя переносить напрямую: юмористическое мышление не всегда является линейным доказательством. Однако сама идея промежуточного поиска оказывается релевантной.

\citet{chen2023prompt} рассматривают генерацию юмора как процесс, который можно направлять пошаговыми инструкциями, основанными на принципах комедийного письма. Такой подход показывает, что явная организация творческого процесса может улучшать ответы по сравнению с простой генерацией по запросу. \citet{Tikhonov2024} развивают похожую идею для коротких шуток, реконструируя процесс их создания как многошаговый механизм. В их постановке модель не просто выдаёт финальную шутку, а проходит через промежуточные этапы, связанные с поиском предпосылки, неожиданного поворота и окончательной формулировки.

Другие работы подчёркивают, что творческие задачи требуют не только последовательного рассуждения, но и ассоциативных скачков. \citet{zhong2024let} предлагают парадигму Leap-of-Thought для генерации юмора в формате Oogiri. В отличие от стандартной цепочки рассуждения, такая схема акцентирует внимание на неожиданных ассоциациях между отдалёнными понятиями. \citet{Kim2025} описывают генерацию юмора как сочетание когнитивных, социальных и творческих навыков, а \citet{ZhangJ2025} предлагают многоэтапную теоретически направляемую схему для интерпретируемой мультимодальной генерации юмора.

Эти исследования важны для данной работы по двум причинам. Во-первых, они подтверждают, что промежуточные представления и творческий поиск могут быть полезны для генерации юмора. Во-вторых, большинство таких подходов задаёт структуру рассуждения вручную: исследователь заранее определяет этапы, их порядок и типы промежуточных объектов. В данной работе используется другая перспектива. Траектория рассуждения рассматривается как случайная переменная, которую модель может порождать сама. Обучение должно не только повысить вероятность конкретных эталонных шуток, но и изменить распределение траекторий так, чтобы они чаще приводили к правдоподобным юмористическим ответам.

\section{Ограничения стандартного постобучения}

Наиболее прямой способ улучшить генерацию юмора состоит в дообучении модели с учителем на корпусе шуток. Однако такая постановка плохо соответствует структуре задачи. Для одного запроса существует много допустимых шуток, и наблюдаемая в корпусе шутка является только одним из возможных вариантов. Если модель обучается воспроизводить единственный ответ, она может улучшать стилистическую похожесть на корпус, но не обязательно улучшать способность к творческому обобщению.

Общая литература по постобучению усиливает это опасение. \citet{Chu2025} противопоставляют дообучение с учителем и обучение с подкреплением, показывая, что первое может сильнее запоминать обучающие шаблоны, тогда как второе при удачной постановке лучше переносится на новые варианты задач. Для юмора это различие особенно важно: запоминание частотных форм шуток не равно способности построить новый неожиданный панчлайн.

Оптимизация по предпочтениям и обучение с подкреплением по модели награды также имеют ограничения. \citet{Gao2023} показывают, что чрезмерная оптимизация несовершенной модели награды может ухудшать истинное качество. В творческих задачах этот риск особенно велик, потому что модель награды или модель-оценщик неизбежно отражает неполные и шумные критерии. Работа \citet{Zhang2024HumorAI} показывает ограничения RLHF- и DPO-подобных методов для юмористических подписей к карикатурам, а результаты \citet{Sakabe2025} и \citet{winters2025evaluating} подчёркивают расхождение между модельными и человеческими оценками юмора.

Следовательно, стандартные методы постобучения не решают задачу полностью. Дообучение с учителем может усиливать имитацию корпуса, а обучение по внешней модели награды может оптимизировать неполный или смещённый критерий. Для данной работы это мотивирует поиск промежуточной постановки: использовать корпус эталонных шуток как слабый сигнал качества, но не обучаться на нём как на единственно правильных ответах и не вводить отдельную модель награды.

\section{Обучение без внешнего верификатора}

Недавние успехи обучения с подкреплением для рассуждающих языковых моделей во многом связаны с доменами, где доступна объективная проверка ответа. В математике и программировании можно задать верификатор, который проверяет правильность результата. \citet{Lightman2023} показывают пользу процессной разметки для математических задач, а \citet{DeepSeek2025} демонстрируют, что обучение с подкреплением может вызывать длинные траектории рассуждения даже без заранее написанных человеком цепочек рассуждения. Однако такие результаты опираются на возможность проверить ответ, что плохо переносится на творческие задачи.

Подходы без внешнего верификатора пытаются заменить явную проверку внутренними вероятностями самой модели. \citet{Zhou2025} предлагают VeriFree: вместо того чтобы генерировать ответ и проверять его внешним модулем, метод максимизирует вероятность эталонного ответа при условии запроса и сгенерированной траектории рассуждения. В такой постановке вероятность эталонного ответа одновременно играет роль сигнала награды для траектории рассуждения и источника прямого градиента для увеличения правдоподобия эталонного ответа.

Близкую идею развивают \citet{Yu2025} в методе RLPR. Они используют внутренние вероятностные оценки модели как сигнал награды и подчёркивают необходимость стабилизации такого сигнала, поскольку вероятности длинных ответов имеют высокую дисперсию и могут становиться численно неудобными. Этот вывод особенно важен для генерации юмора. В задачах с коротким ответом вероятность эталонной последовательности ещё может быть практическим сигналом. Для полной шутки, состоящей из десятков токенов, произведение токеновых вероятностей быстро становится очень малым и начинает сильно зависеть от длины ответа.

Ограничение методов VeriFree и RLPR для данной работы состоит в том, что они в первую очередь формулируются вокруг одного эталонного ответа или вокруг задач, где эталонный ответ можно считать достаточно определённым. Юмор устроен иначе. Один запрос может иметь множество хороших ответов, а наблюдаемое множество шуток является неполным и смещённым подмножеством возможных удачных вариантов. Поэтому прямое применение одноэталонной цели может ошибочно наказывать хорошие, но не наблюдавшиеся шутки, и чрезмерно усиливать отдельные формулировки из корпуса.

\section{Позиционирование данной работы}

Данная работа находится на пересечении трёх направлений: генерации юмора, обучения рассуждающих языковых моделей и обучения без внешнего верификатора. Литература по вычислительному юмору показывает, что задача является многомерной, контекстно зависимой и плохо сводится к единственной метрике \citep{ritchie2001current,amin2020survey,Loakman2025,lemmens2026computational}. Работы по генерации юмора с рассуждением показывают, что промежуточный творческий поиск может улучшать ответы, но часто задают структуру этого поиска вручную \citep{chen2023prompt,zhong2024let,Tikhonov2024,Kim2025,ZhangJ2025}. Работы по обучению без внешнего верификатора показывают, что вероятность эталонного ответа может служить внутренним сигналом обучения, но их исходная постановка плохо учитывает многореференсность и неполноту творческих задач \citep{Zhou2025,Yu2025}.

Основной исследовательский пробел состоит в том, что существующие методы либо рассматривают юмор как задачу генерации по вручную заданному сценарию рассуждения, либо рассматривают обучение рассуждений в доменах с объективной проверкой или одним эталонным ответом. В данной работе предлагается адаптировать идею обучения без внешнего верификатора к юмору как к частично наблюдаемой многореференсной задаче. Вместо единственного эталонного ответа используется множество эталонных шуток, найденных в корпусе по близости запроса. Вместо сырой вероятности полной шутки используется нормированная логарифмическая масса эталонных ответов, что уменьшает численные проблемы и смещение в пользу коротких последовательностей. Траектория рассуждения при этом остаётся порождаемой моделью, а не задаётся фиксированным шаблоном.

Такое позиционирование определяет вклад работы. Она не утверждает, что автоматическая оценка юмора полностью заменяет человеческую, и не предполагает, что наблюдаемые эталонные шутки исчерпывают все хорошие ответы. Напротив, работа рассматривает эти ограничения как часть постановки задачи. Цель состоит в разработке и экспериментальной проверке программного конвейера, который позволяет изучать, можно ли использовать многореференсную логарифмическую цель для обучения рассуждающей генерации юмора без отдельной модели награды.